% META: {"id": "1NSI-TYPES-BASE-RE-C5", "chapitre": "1NSI-TYPES-BASE", "type_objet": "remediation", "capacites": ["P-BASE-05"], "capacites_codes": ["C5"], "mode_creation": "ex_nihilo", "sources_inspiration": [], "status": "needs_review"}

%---------------------------------------------------------------
% FICHE DE REMEDIATION — C5 : plusieurs encodages conviennent souvent
%---------------------------------------------------------------

\textbf{Ce qui bloque.} Tu retiens qu'UTF-8 représente tous les caractères — c'est exact — et
tu en conclus qu'il est le \emph{seul} à convenir. Ou l'inverse : tu choisis ISO-8859-1 en
oubliant qu'ASCII et UTF-8 conviennent aussi dès que le texte n'a pas d'accent.

\textbf{À retenir.} La question n'est jamais « quel encodage est le meilleur » mais « lesquels
peuvent représenter \emph{ce} texte ». Les $128$ premiers caractères d'UTF-8 sont exactement
ceux d'ASCII : un texte sans accent s'encode à l'identique dans les trois.

\begin{exercice}{1NSI-TB-RE-C5-EX1}{1}{10}
Pour chacune des trois chaînes \lstinline{"Note: 12/20"}, \lstinline{"Élève"} et
« \lstinline{prix 15} » suivi du symbole euro :
\begin{enumerate}
  \item Indiquer lesquels d'ASCII, ISO-8859-1 et UTF-8 peuvent la représenter.
  \item Donner sa taille en octets en UTF-8, et la comparer au nombre de caractères.
  \item Expliquer pourquoi le symbole euro se comporte différemment du \lstinline{É}.
\end{enumerate}
\end{exercice}

% BEGIN-VERIFY
% assert len("Note: 12/20") == 11
% assert len("Note: 12/20".encode("ascii")) == 11
% assert len("Note: 12/20".encode("utf-8")) == 11
% assert len("Élève") == 5
% assert len("Élève".encode("utf-8")) == 7
% assert len("Élève".encode("iso-8859-1")) == 5
% assert len("prix 15 €") == 9
% assert len("prix 15 €".encode("utf-8")) == 11
% try:
%     "prix 15 €".encode("iso-8859-1")
%     euro_dans_iso = True
% except UnicodeEncodeError:
%     euro_dans_iso = False
% assert euro_dans_iso is False
% END-VERIFY
